\chapter{Trabalhos Relacionados} \label{chapt:trabalhos}

O \ac{JITJSS} foi proposto pela primeira vez em 2008 por \cite{baptiste2008lagrangian}. Desde então, diferentes autores estudaram o problema e propuseram diferentes estratégias para resolvê-lo. Nesta seção, listaremos com detalhes os trabalhos que abordaram o problema. Além disso, ainda não há nenhum trabalho que aplicou as técnicas de \ac{AAD} especificamente ao \ac{JITJSS}. Por isso, como forma de justificar este trabalho, apresentaremos os usos de \ac{AAC} e \ac{AAD} da literatura para encontrar algoritmos estado da arte para diversos problemas. 
 
\section{\textit{Just-In-Time Job Shop}} \label{sect:bibjitjsp}

\cite{baptiste2008lagrangian} propuseram as 72 instâncias do problema que foram utilizadas como métrica de qualidade em todos os outros trabalhos que estudaram o \ac{JITJSS}. Os autores também desenvolveram duas relaxações lagrangianas para encontrar limites inferiores para cada instância. Essas estratégias relaxam as restrições de precedência das tarefas e as restrições de máquina. A partir desses limites inferiores, eles utilizam métodos de \textit{branch-and-bound} e programação linear para encontrar soluções viáveis e uma busca local para melhorar as soluções encontradas. Nesse sentido, o trabalho propôs os primeiros limites superiores para as instâncias, mas, em diversas instâncias, a diferença entre os limites inferiores e superiores ainda era grande.

\cite{laborie2007self} desenvolveram um método genérico para resolver problemas de agendamento que utiliza uma busca auto-adaptativa em grandes vizinhanças, baseada em relaxar soluções através de um portfólio de vizinhanças e reotimizá-las. Com essa abordagem, eles conseguiram limites superiores melhores em diversas instâncias. \cite{monette2009just} também utilizaram um processo de relaxação das máquinas, porém, aplicado em um algoritmo \textit{branch-and-bound} desenvolvido por eles. Os autores também acoplaram dois métodos de busca local no processo e, com isso, conseguiram melhorar o estado da arte do problema.

Outro tipo de abordagem utilizada são os algoritmos evolucionários. \cite{araujo2009genetic} propuseram um algoritmo genético combinado com uma busca local que foi expandido por \cite{dos2010combination}. No segundo trabalho, os autores utilizam um método de programação linear para gerar um agendamento através de uma sequência. Isso torna o modelo menos flexível e mais contido, gerando agendamentos ótimos em menos tempo. Com essas abordagens, os autores melhoraram, na época, 56 das 72 instâncias do estado da arte.

Ainda no âmbito de algoritmos evolucionários, \cite{yang2012improved} propõem em seu trabalho um algoritmo genético que difere dos anteriores, de forma mais significativa, na decodificação de um cromossomo. Os autores desenvolveram um método dividido em três estágios para gerar agendamentos válidos. Esse método obteve resultados ligeiramente melhores em algumas instâncias de 10 e 15 tarefas, e os autores não reportaram os resultados para as instâncias maiores.

A meta-heurística \ac{VNS} também já foi aplicada ao problema. Em \cite{wang2014variable}, os autores propõem um \ac{VNS} com duas vizinhanças. Além disso, eles utilizam a mesma estratégia utilizada por \cite{dos2010combination} na geração de agendamentos, isto é, utilizam uma meta-heurística para explorar o espaço de sequências e um método de programação inteira para gerar um agendamento ótimo para cada sequência. Os autores conseguiram encontrar vários resultados melhores que o estado da arte, principalmente nas instâncias de 10 e 15 tarefas. 

\cite{ahmadian2021meta} também utilizaram um \ac{VNS} para resolver o problema. Porém, nesse trabalho, além de eles utilizarem programação inteira para gerar agendamentos a partir de sequências, eles utilizam esses métodos para melhorar as próprias sequências. Para isso, os autores propuseram selecionar um subconjunto das tarefas de uma máquina e relaxar as restrições de precedência dessas operações. Por conseguinte, o modelo passa a procurar a ordem ótima apenas para esse subconjunto de operações selecionadas. Essa estratégia simplifica o modelo de programação inteira, gerando resultados melhores em menos tempo. Essa proposta conseguiu atingir ou melhorar o estado da arte em 61 das 72 instâncias, tornando-se a solução mais eficiente e abrangente da literatura. 

Mais recentemente, \cite{abderrazzak2022adaptive} propuseram novamente uma busca em grandes vizinhanças para resolver o problema. Essa abordagem, porém, consiste em destruir e reconstruir as soluções de forma adaptativa durante a execução do algoritmo. Os autores conseguiram atingir o estado da arte na maior parte das instâncias de 10 e 15 tarefas e melhorar algumas instâncias pontualmente. Além dessa estratégia, \cite{sabri2023reinforcement} recentemente propuseram uma abordagem ao problema com aprendizado por reforço profundo que obteve poucas melhorias na qualidade das soluções, mas obteve soluções competitivas para as instâncias de 10 e 15 tarefas utilizando pouco tempo de processamento em comparação com o restante da literatura. Ambos os trabalhos não reportaram os resultados para as maiores instâncias do problema.

Ao longo dos anos, esses trabalhos encontraram soluções cada vez melhores para as instâncias propostas por \cite{baptiste2008lagrangian}. Nesse sentido, as melhores soluções encontradas para cada instância estão contidas em cinco trabalhos: \cite{dos2010combination, wang2014variable, ahmadian2021meta, abderrazzak2022adaptive, sabri2023reinforcement}. Dessa forma, quando nos referirmos ao estado da arte para o problema \ac{JITJSS}, estaremos referenciando a composição dos resultados desses cinco trabalhos.

\section{Projeto Automático de Algoritmos} \label{sect:bibconfig}

Em \cite{liao2014unified}, os autores desenvolveram um Otimização por Colônia de Formigas
(do inglês, Ant Colony Optimazation -- \acs{ACO}) para otimização contínua que combina componentes de outros métodos \ac{ACO} da literatura em um ambiente genérico de configuração. Dessa forma, o algoritmo desenvolvido pode instanciar tanto as estratégias da literatura quanto combinações de componentes nunca vistas. Os autores utilizaram configuração automática para instanciar os algoritmos e, com isso, eles obtiveram resultados melhores em comparação a outros trabalhos da literatura.

Os trabalhos de \cite{aydin2017abc} e \cite{bezerra2020automatically} realizam um processo parecido no sentido de reunir estratégias verificadas da literatura para uma meta-heurística específica, mas eles abordam, respectivamente, os algoritmos de colônia artificial de abelhas e algoritmos evolucionários multiobjetivos. Em ambos os trabalhos, os autores reportaram que a performance dos algoritmos gerados era maior quando comparada a outros algoritmos da literatura.

Com uma abordagem menos centrada no tipo de meta-heurística utilizada, \cite{khudabukhsh2016satenstein} utilizam uma abordagem \textit{top-down} de projeto automático de algoritmos para gerar buscas locais estocásticas que resolvem o problema da \ac{SAT}. Eles comparam os algoritmos gerados pela sua abordagem com outros 11 \textit{solvers} do \ac{SAT} e reportam resultados melhores do que os obtidos pelos outros \textit{solvers}. 

Já \cite{de2018automatic} utilizam uma abordagem \textit{bottom-up} para gerar heurísticas automaticamente para a \ac{UBQP}. Os autores reportam resultados competitivos com estado da arte obtidos pelos algoritmos gerados. Em \cite{de2018automatically}, os autores apresentam uma redução do problema de alocação de provas para \ac{UBQP} e utilizam novamente uma abordagem \textit{bottom-up} para gerar heurísticas para o problema. A abordagem proposta obteve resultados melhores que o estado da arte para o problema de alocação de provas, mostrando-se uma alternativa para resolver também outros problemas redutíveis ao \ac{UBQP}.

No âmbito de problemas de agendamento, as variações do \ac{FSP} foram muito abordadas com técnicas de projeto automático. \cite{marmion2013automatic} propõem a estrutura de uma busca local estocástica híbrida genérica para gerar algoritmos para o problema \ac{PFSP-WT}. Os autores geraram três algoritmos diferentes. Com isso, eles compararam os resultados com o algoritmo estado da arte do problema e, para cada instância, pelo menos um dos algoritmos era melhor que o estado da arte, sendo que, na maioria das instâncias, os três algoritmos gerados eram melhores que o estado da arte.

\cite{pagnozzi2019automatic} também propuseram um espaço de projeto automático de buscas locais estocásticas híbridas. Nesse trabalho, os autores abordaram três funções objetivo diferentes para o \ac{PFSP}. Em comparação com o estado da arte para cada tipo de função objetivo, dois dos três algoritmos gerados se tornaram o novo estado da arte para as suas respectivas funções objetivo.

\cite{brum2018automatic} também abordam o \ac{PFSP} através do projeto automático de algoritmos. Porém, neste trabalho, os autores desenvolvem um modelo de geração de algoritmos baseado na meta-heurística \ac{ILS}. Os autores geraram 10 algoritmos e compararam a performance deles contra o estado da arte. Dessa maneira, eles obtiveram resultados melhores que o estado da arte com a maioria dos algoritmos gerados.

\cite{alfaro2020automatic} propõem soluções para três funções objetivo do problema \ac{HFS}. Eles também utilizam um modelo de \ac{ILS} para gerar algoritmos automaticamente. Os autores conseguiram gerar dois novos algoritmos estado da arte para as funções objetivo \ac{TFT} e Atraso e Adiantamento Totais Ponderados.

\cite{brum2022automatic} aborda outro problema chamado \ac{NPFSP}. Os autores dividem o processo em duas fases. A primeira consiste em resolver a instância com as restrições do \ac{PFSP}. A segunda fase consiste em resolver o \ac{NPFSP} utilizando a solução da primeira fase como solução inicial. A partir dessa ideia, os autores geram automaticamente algoritmos para o \ac{NPFSP}. Os algoritmos gerados obtiveram melhores resultados que o estado da arte tanto para o \ac{PFSP} quanto para o \ac{NPFSP}; isto é, os resultados gerados pela primeira fase já eram melhores que o estado da arte para o \ac{PFSP} e, ao receber essas soluções como entrada, a segunda fase também obteve resultados melhores que o estado da arte para o \ac{NPFSP}.

Todos esses trabalhos têm como característica comum obter algoritmos que superaram total ou parcialmente o estado da arte em seus respectivos domínios com as técnicas de \ac{AAC} e \ac{AAD}. Além disso, apesar de nenhum trabalho ter aplicado essas técnicas ao \ac{JITJSS}, temos diversos trabalhos que demonstraram a eficácia da \ac{AAC} e do \ac{AAD} em outros problemas de agendamento. A partir disso, o presente trabalho se sustenta na hipótese de que a aplicação das técnicas de \ac{AAC} e \ac{AAD} ao \ac{JITJSS} resultará em novos algoritmos estado da arte para o problema. 
